\chapter{DATA DICTIONARY AND SCORING RULES}
\label{app:data}
\newpage
\section{Record Identity and Provenance}
This dictionary defines the meaning of each recorded field. Names are descriptive labels, not exporter field names, and will be mapped to the versioned exporter schema before data collection. Some listed fields may not yet be implemented in the current build.

\begingroup
\thesislongtablecaption{Record identity and scope fields.}
\begin{longtable}{@{}p{0.26\linewidth}p{0.65\linewidth}@{}}
\toprule
Field family & Meaning \\\midrule\endfirsthead
\toprule Field family & Meaning \\\midrule\endhead
Participant ID & Coded study identifier; not a name or voice-provider credential.\\
Session and run & Unique attempt/session context, including development versus study mode.\\
Block and trial & Referenced measured block, within-block trial, global trial index, and separate practice flag.\\
Voice condition & Neutral or self-similar; actual voice ID, provider, synthesis model and configuration.\\
Assigned order & Neutral-first or self-similar-first; practice uses the first assigned voice.\\
Build and policy & Application revision, configuration, state-machine policy, phrase semantics, and library version.\\
Schedule & Seed, schedule version, nominal wall assignment, direction and height assignment, realized geometry, and attempt status.\\
Survey scope & Administration phase and timestamp, recorded separately from the referenced participant, block, and voice; each side-by-side rating carries its voice column (Selected Voice or Self-Similar Voice).\\
Intake & Audio intake answers (eligibility and background), the intake agent voice ID, and the identifiers of the speech samples used for cloning, linked to the participant's study ID.\\\bottomrule
\end{longtable}
\endgroup

\section{Geometry and Timing Fields}
\begingroup
\thesislongtablecaption{Geometry and event fields.}
\begin{longtable}{@{}p{0.28\linewidth}p{0.18\linewidth}p{0.43\linewidth}@{}}
\toprule
Quantity & Unit/reference & Definition \\\midrule\endfirsthead
\toprule Quantity & Unit/reference & Definition \\\midrule\endhead
Nominal angle & Degrees; wall schedule & Allocation category, distinct from exact target bearing.\\
Exact starting theta & Degrees; pre-cue cross & Signed horizontal angle from saved origin-to-cross toward origin-to-target; saved at readiness completion; absent from earlier recordings.\\
Height difference & Metres; cross center & Target Y minus cross Y; intended and realized sign stored separately.\\
Replay wall bounds & Metres; recording-local frame & Wall bottom, height, and availability flag, saved in current rotational-search snapshots. For earlier recordings, bounds are estimated from objects and saved positions are not rescaled.\\
Target distance & Metres; first tracked head after reveal & Three-dimensional distance with observation timestamp.\\
Reveal & Session clock & Object visibility and active-search onset event.\\
First gaze contact & Active clock & First valid target collider intersection; missing if absent.\\
Last gaze contact & Active clock & Last valid observed target intersection; a contact time, not a fixation duration.\\
First controller contact & Active clock & First tracked right-ray intersection with target.\\
Gaze qualification & Session clock & Time at which valid winning-vicinity evidence satisfies the interface timer.\\
Pointing qualification & Session clock & Time at which controller-hold and current-gaze requirements are met.\\
Correct selection & Session and active clocks & Trigger selection of target, recorded separately from acknowledgment.\\
Pause intervals & Session clock & Start, end, and reason; their union is used to compute active time.\\
Instruction onset & Session clock & Actual software playback start, phrase and event category.\\\bottomrule
\end{longtable}
\endgroup

\section{Raw Observation Records}
Gaze records will retain available source time or sequence, application time, validity, source identity, origin, direction, and sample age. Per-eye data will be labeled as such when available, and combined gaze will not be treated as independent left- and right-eye observations. Any head-direction fallback will be flagged rather than silently substituted for eye gaze in eye-derived outcomes.

Motion records will retain head and right-controller poses, tracking state, trigger state, and the time basis needed for reconstruction. Object records will retain stable identifiers, shape/color metadata, transforms, target designation, wall assignment, and layout version. Event records will identify state transitions, eligible actions, model requests and responses, cancellations, playback, interruptions, and trial outcomes. Unknown values are marked unavailable rather than filled with plausible defaults.

\section{Scoring Rules}
NASA-TLX~\cite{hart1988nasatlx} scoring retains the six raw values on the 0--10 scale. The positively phrased Performance item runs from 0 = Failure to 10 = Perfect, the reverse of the NASA form, so larger raw values indicate better perceived performance; demand, effort and frustration keep their adverse direction. The raw Performance response is preserved, and a separate workload-oriented value is calculated once as $10-r_{\mathrm{Performance}}$. The unweighted workload mean combines this reversed Performance value with the other five raw ratings, none of which is reversed, and is calculated only when all six ratings are present; higher means indicate greater workload or poorer perceived performance. Any exported normalized 0--100 field is named separately, calculated from explicit bounds, and labeled as raw or workload-oriented. Legacy data without known bounds or Performance direction remain unscored, and earlier Perfect-to-Failure responses are not reversed automatically.

IMI~\cite{ryan1982control,mcauley1989imi} scoring uses separate complete subscale means: interest/enjoyment (7), effort/importance (5), pressure/tension (5), and value/usefulness (7). On the 1--7 response scale, reverse-coded values use $8-r$. Within-subscale reverse positions are enjoyment 3 and 4, effort 2 and 5, pressure 1 and 3, and none for value. There is no overall IMI total.

The agent-perception battery selected from Guo et al.~\cite{guo2024doppelganger} retains source identifiers. Co-presence Q1--6, perceived intelligence Q13--18, and anthropomorphism Q37--42 are separate six-item means when complete. Q20 appealing, Q21 eerie, and Q22 familiar are individual ratings in their original directions. There is no 21-item or combined eeriness total. Custom voice/assistance items likewise remain individual outcomes. The five draft ARQ items~\cite{wang2020terrible} (helpfulness, relatability, appropriate behavior, likeability, willingness to interact again) are individual 1--7 ratings with no composite score.

\section{Export Validation}
Before analysis, automated checks will reject duplicate trial keys, unknown voice/block joins, negative active intervals, selection before reveal, and scores outside declared bounds, while preserving incomplete outcomes, missing contacts, absent questionnaire items, and old records with insufficient provenance. Joined tables will be checked for exactly one survey observation per participant and referenced block before survey values are repeated on trial rows. Exports will be reproducible from raw records and a versioned analysis configuration.
